%%%%%%%%%%%%%%%%%%%%%%%%%%%%%%%%%%%%%%%%%%%%%%%%%%%%%%%%%%%%%%%%%%%%%%%%%%%%%%%%%%%%
%Do not alter this block of commands.  If you're proficient at LaTeX, you may include additional packages, create macros, etc. immediately below this block of commands, but make sure to NOT alter the header, margin, and comment settings here. 
\documentclass[12pt]{article}
\usepackage{babel}
\usepackage[margin=1in]{geometry} 
\usepackage{amsmath,amsthm,amssymb,amsfonts, enumitem, fancyhdr, color, comment, graphicx, environ}
\pagestyle{fancy}
\setlength{\headheight}{65pt}
\newenvironment{problem}[2][Question] %{\begin{trivlist}
%\item[\hskip \labelsep {\bfseries #1}\hskip \labelsep {\bfseries #2.}]}{\end{trivlist}}
\newenvironment{sol}
    {
    }
    {
    }
\specialcomment{com}{ \color{blue} \textbf{Comment:} }{\color{black}} %for instructor comments while grading
\NewEnviron{probscore}{\marginpar{ \color{blue} \tiny Problem Score: \BODY \color{black} }}
%%%%%%%%%%%%%%%%%%%%%%%%%%%%%%%%%%%%%%%%%%%%%%%%%%%%%%%%%%%%%%%%%%%%%%%%%%%%%%%%%
\usepackage{enumitem}
\usepackage[linguistics]{forest}
\usepackage{ragged2e}
\usepackage{nccmath, mathtools}
\newcommand\mytab[1][1cm]{\hspace*{#1}}
\usepackage{tabto}
\usepackage{fontspec}
\usepackage{gb4e}
% \babelfont[tibetan]{rm}
%             {Jomolhari}

%%%%%%%
% section number modifier 
\newcommand{\mysection}[2]{\setcounter{section}{#1}\addtocounter{section}{-1}\section{#2}}


%%%%%%%%%%%%%%%%%%%%%%%%%%%%%%%%%%%%%%%%%%%%%
%Fill in the appropriate information below
\lhead{Arav Sarma \\ 15193845}  %replace with your name
\rhead{Homework 0 \\ 2025 Speech Processing} %replace XYZ with the homework course number, semester (e.g. ``Spring 2019"), and assignment number.
%%%%%%%%%%%%%%%%%%%%%%%%%%%%%%%%%%%%%%%%%%%%%


%%%%%%%%%%%%%%%%%%%%%%%%%%%%%%%%%%%%%%
%Do not alter this block.
\begin{document}

%%%%%%%%%%%%%%%%%%%%%%%%%%%%%%%%%%%%%%

\mysection{0}{Reflection: How processing works}
%\section{\large\textbf{Reflection: How processing works}}


\subsection*{\emph{a.}}
Without influence of absense/presence of a word, information will move up unidirectionally; auditory information will be categorized into phonological structure without top-to-bottom interference from lexical knowledge, from here if the information retrieved exists in phono\-logical storage (or directly as morpheme) of lexicon, that lexical form will be accessed, if not, no lexical access will occur (unless other mechanism is present which approximates possible lexical forms from the categorized phonological forms without needing any auditory information).

\subsection*{\emph{b.}}
With influence from presence/absence of a word, categorisation of auditory information will have additional interference from lexical knowledge, this means information is processed both (simultaneously) bottom-top and top-bottom (there may be various mechanisms for this; goodness of fit, echoic store) and contexts for when such mechanisms arrise/are needed. The context Ganong looked at is ambiguous candidates near the phonemic boundary on VOT continua. Here lexical access occurs for ambiguous candidates when one side of the continuum is an existing morpheme/word, therefore being susceptible to top-bottom interference, whereas candidates from extreme ends do not show this effect. 

\subsection*{\emph{c.}}
Option 2


% \section{\large\textbf{Payne 8.3: Avar}}

% \subsubsection*{\emph{Question A}}

% \textbf{Language family:} Nakh-Daghestanian (→ Daghestanian→ Avar-Andic-Tsezic)

% \subsubsection*{\emph{Question B}}

% \begin{exe}
%     \exi{(6)}
%     \gll  Yas-as vas v-etts-ana. \\
%     Girl-\textsc{ERG} boy.\textsc{ABS} \textsc{ABS.M(?)}-praise-\textsc{PST} \\
%     \glt  `The girl praised the boy.' \\
% \end{exe}

% \subsubsection*{\emph{Question D}}
% \justifying
% \textbf{Word Order:} \par
% \vspace*{2mm}
% Neutral alignment, A, P and S, all are post verbal, not clear whether A and P are neutralised, as they are not neutralised in case it may allow for both AP and PA orders, dataset in the question only gives AP; \par
% APV: 5. \emph{Vasas} (A) \emph{yas} (P) \emph{yettsana} (V), \par
% SV: 4. \emph{Yas} (S) \emph{yegana} (V) \\
% \\
% \justifying
% \vspace*{2mm}
% \textbf{Case:} \par
% \vspace{2mm}
% Ergative alignment, S = P vs A. S and P are unmarked, A is marked with suffix \emph{\textbf{-as}}. \par
% 5. \emph{Vas-\textbf{as}} (ERG) \emph{yas-\textbf{Ø}} (ABS) \emph{yettsana}, \par
% 4. \emph{Yas-\textbf{Ø}} (ABS) \emph{yegana} \vspace*{4mm}
% \\
% \justifying
% \vspace*{2mm}
% \textbf{Agreement:} \par
% \vspace*{2mm}
% Ergative alignment, S = P vs A. S and P are marked with prefix \emph{\textbf{y-/v-}}, A is unmarked. \par
% 5. \emph{Vasas} (A) \emph{yas} (P) (ABS)\emph{\textbf{y-}}\emph{ettsana-\textbf{Ø}}(ERG), \par
% 4. \emph{Yas} (S) (ABS)\emph{\textbf{y-}}\emph{egana}


% \section{\large\textbf{Language Design}}

% \subsubsection*{\emph{Question 1: Language with active alignment of agreement}}
% \justifying

% \paragraph{Intransitive}
% \begin{exe}
%     \exi{(S_{a})}
%     \gll  Kayap-Ø mekan-nung-Ø \\
%     hunter-\textsc{ABS} swim-\textsc{PST}-\textsc{A} \\
%     \glt  `The/a hunter swam.' \\
% \end{exe}
% \begin{exe}
%     \exi{(S_{p})}
%     \gll  Kayap-Ø lacat-tung-ngiq \\
%     hunter-\textsc{ABS} die-\textsc{PST}-\textsc{3.P.AN} \\
%     \glt  `The/a hunter died.' \\
%     AN = animate
% \end{exe}
% \paragraph{Transitive}
% \begin{exe}
%     \exi{(A,P)}
%     \gll  Kayap-pit tyupin-Ø quy-yung-ngip-Ø \\
%     hunter-\textsc{ERG} pigeon-\textsc{ABS} eat-\textsc{PST}-\textsc{3.P.INAN}-\textsc{A} \\
%     \glt  `The/a hunter ate the/a pigeon.' \\
%     INAN = inanimate
% \end{exe}

% \subsubsection*{\emph{Question 2: Verb initial word order, split agreement conditioned by person (3rd: neutral, 1st,2nd: accusative)}}
% \justifying

% \paragraph{3rd person}
% \begin{exe}
%     \exi{(S_{3})}
%     \gll  sx'-kts-Ø tlqwt \\
%     \textsc{PST}-sleep-3.S hunter \\
%     \glt  `The/a hunter slept.'
% \end{exe}
% \begin{exe}
%     \exi{(A_{3},P_{3})}
%     \gll  sx'-lsk'-Ø-Ø tlqwt yxtš \\
%     \textsc{PST}-kill-3.A-3.P hunter pigeon \\
%     \glt  `The/a hunter killed the/a pigeon.'
% \end{exe}
% \paragraph{1st-2nd person}
% \begin{exe}
%     \exi{(S_{1})}
%     \gll  sx'-kts-sw sqw \\
%     \textsc{PST}-sleep-1.NOM 1 \\
%     \glt  `I/we slept.'
% \end{exe}
% \begin{exe}
%     \exi{(S_{2})}
%     \gll  sx'-kts-w'x kw'ts \\
%     \textsc{PST}-sleep-2.NOM 2.SG \\
%     \glt  `You slept.'
% \end{exe}
% \begin{exe}
%     \exi{(A_{1},P_{2})}
%     \gll  sx'-lsk'-sw-Ø sqw kw'ts \\
%     \textsc{PST}-kill-1.NOM-2.ACC 1 2.SG \\
%     \glt  `I/we killed you.'
% \end{exe}
% \begin{exe}
%     \exi{(A_{2},P_{1})}
%     \gll  sx'-lsk'-w'x-Ø kw'ts sqw \\
%     \textsc{PST}-kill-2.NOM-1.ACC 2.SG 1 \\
%     \glt  `You killed me/us.'
% \end{exe}

% This system fits with typological generalisations, as it is common for V-initial languages to have non-neutral alignment for agreement and not necessarily for case. Although it does not have overt marking for both arguments, which is more common for agreement strategy, the argument with agreement is A which is more preferred than having agreement with only P. Typologically it is also typical to have split with 1st-2nd vs 3rd, where by following the referential hierarchy, 1st-2nd (leftmost in hierarchy) are accusatively aligned and 3rd is not (neutrally aligned). 

%%%%%%%%%%%%%%%%%%%%%%%%%%%%%%%%%%%%%
%Do not alter anything below this line.
\end{document}